% ============================================================
%  EdgeLeak AI — Guía de Integración por API
%  Documento técnico para equipo externo (Flutter + Firebase)
%  Versión 1.0 — 2025
% ============================================================
\documentclass[12pt, a4paper]{article}

% --- Paquetes ---
\usepackage[utf8]{inputenc}
\usepackage[T1]{fontenc}
\usepackage[spanish]{babel}
\usepackage{geometry}
\usepackage{graphicx}
\usepackage{xcolor}
\usepackage{listings}
\usepackage{hyperref}
\usepackage{booktabs}
\usepackage{array}
\usepackage{fancyhdr}
\usepackage{titlesec}
\usepackage{enumitem}
\usepackage{mdframed}
\usepackage{fontawesome5}
\usepackage{lmodern}
\usepackage{microtype}
\usepackage{parskip}
\usepackage{tcolorbox}
\tcbuselibrary{skins, breakable}

% --- Márgenes ---
\geometry{
  top=2.5cm, bottom=2.5cm,
  left=2.8cm, right=2.8cm
}

% --- Colores corporativos EdgeLeak ---
\definecolor{edgeblue}{RGB}{13, 71, 161}
\definecolor{edgelightblue}{RGB}{227, 242, 253}
\definecolor{edgegray}{RGB}{96, 125, 139}
\definecolor{edgegreen}{RGB}{27, 94, 32}
\definecolor{edgered}{RGB}{183, 28, 28}
\definecolor{edgeorange}{RGB}{230, 81, 0}
\definecolor{codebg}{RGB}{245, 245, 245}
\definecolor{codeframe}{RGB}{200, 200, 200}
\definecolor{successgreen}{RGB}{46, 125, 50}
\definecolor{warningyellow}{RGB}{245, 124, 0}

% --- Configuración de hipervínculos ---
\hypersetup{
  colorlinks=true,
  linkcolor=edgeblue,
  urlcolor=edgeblue,
  pdftitle={EdgeLeak AI - Guía de Integración por API},
  pdfauthor={Equipo EdgeLeak AI},
}

% --- Cabecera y pie de página ---
\pagestyle{fancy}
\fancyhf{}
\fancyhead[L]{\textcolor{edgegray}{\small EdgeLeak AI}}
\fancyhead[R]{\textcolor{edgegray}{\small Guía de Integración por API}}
\fancyfoot[C]{\textcolor{edgegray}{\thepage}}
\renewcommand{\headrulewidth}{0.4pt}
\renewcommand{\footrulewidth}{0pt}

% --- Estilo de secciones ---
\titleformat{\section}
  {\Large\bfseries\color{edgeblue}}
  {\thesection.}{0.5em}{}[\color{edgeblue}\titlerule]
\titleformat{\subsection}
  {\large\bfseries\color{edgeblue!80!black}}
  {\thesubsection.}{0.5em}{}
\titleformat{\subsubsection}
  {\normalsize\bfseries\color{edgegray}}
  {\thesubsubsection.}{0.5em}{}

% --- Estilo de código (JSON / Dart / HTTP) ---
\lstdefinestyle{jsonStyle}{
  language=,
  backgroundcolor=\color{codebg},
  frame=single,
  framerule=0.6pt,
  rulecolor=\color{codeframe},
  basicstyle=\ttfamily\footnotesize,
  breaklines=true,
  breakatwhitespace=false,
  showstringspaces=false,
  keywordstyle=\color{edgeblue}\bfseries,
  stringstyle=\color{edgegreen},
  commentstyle=\color{edgegray}\itshape,
  numbers=none,
  xleftmargin=10pt,
  xrightmargin=10pt,
  tabsize=2,
  captionpos=b,
}

\lstdefinestyle{dartStyle}{
  language=Java,
  backgroundcolor=\color{codebg},
  frame=single,
  framerule=0.6pt,
  rulecolor=\color{codeframe},
  basicstyle=\ttfamily\footnotesize,
  breaklines=true,
  breakatwhitespace=false,
  showstringspaces=false,
  keywordstyle=\color{edgeblue}\bfseries,
  stringstyle=\color{edgegreen},
  commentstyle=\color{edgegray}\itshape,
  numbers=left,
  numberstyle=\tiny\color{edgegray},
  numbersep=8pt,
  xleftmargin=20pt,
  xrightmargin=10pt,
  tabsize=2,
  captionpos=b,
  morekeywords={String, int, double, bool, List, Map, Future, void, async, await,
                final, const, class, return, if, else, try, catch, throw, new,
                import, null, true, false, var},
}

\lstdefinestyle{httpStyle}{
  language=,
  backgroundcolor=\color{edgelightblue},
  frame=single,
  framerule=0.8pt,
  rulecolor=\color{edgeblue!50},
  basicstyle=\ttfamily\small\bfseries,
  breaklines=true,
  showstringspaces=false,
  numbers=none,
  xleftmargin=10pt,
  xrightmargin=10pt,
}

% --- Cajas informativas ---
\tcbset{
  infobox/.style={
    enhanced, breakable,
    colback=edgelightblue,
    colframe=edgeblue,
    fonttitle=\bfseries\small,
    title={\faInfoCircle\ Información},
    left=6pt, right=6pt, top=4pt, bottom=4pt,
  },
  warningbox/.style={
    enhanced, breakable,
    colback=yellow!10,
    colframe=warningyellow,
    fonttitle=\bfseries\small\color{warningyellow},
    title={\faExclamationTriangle\ Importante},
    left=6pt, right=6pt, top=4pt, bottom=4pt,
  },
  successbox/.style={
    enhanced, breakable,
    colback=green!8,
    colframe=successgreen,
    fonttitle=\bfseries\small\color{successgreen},
    title={\faCheckCircle\ Ejemplo exitoso},
    left=6pt, right=6pt, top=4pt, bottom=4pt,
  },
  errorbox/.style={
    enhanced, breakable,
    colback=red!8,
    colframe=edgered,
    fonttitle=\bfseries\small\color{edgered},
    title={\faTimesCircle\ Error},
    left=6pt, right=6pt, top=4pt, bottom=4pt,
  },
}

% ============================================================
\begin{document}
% ============================================================

% --- Portada ---
\begin{titlepage}
  \centering
  \vspace*{2cm}

  {\color{edgeblue}\rule{\linewidth}{2pt}}
  \vspace{0.5cm}

  {\Huge\bfseries\color{edgeblue} EdgeLeak AI}\\[0.4cm]
  {\Large\color{edgegray} Sistema de Detección de Microfugas con IA}\\[1cm]

  {\color{edgeblue}\rule{\linewidth}{1pt}}
  \vspace{1cm}

  {\LARGE\bfseries Guía de Integración por API}\\[0.3cm]
  {\large\color{edgegray} Para equipos externos — Flutter + Firebase}\\[2cm]

  \begin{tcolorbox}[
    enhanced,
    colback=edgelightblue,
    colframe=edgeblue,
    width=0.75\textwidth,
    center,
    left=12pt, right=12pt, top=8pt, bottom=8pt,
  ]
    \centering
    \textbf{Versión:} 1.0\\
    \textbf{Puerto base:} \texttt{8080}\\
    \textbf{Protocolo:} HTTP/1.1\\
    \textbf{Formato de datos:} JSON
  \end{tcolorbox}

  \vspace{2cm}

  {\color{edgegray}\large
    Corporación Universitaria Remington\\
    Proyecto Integrador (PICUR)\\
    \vspace{0.3cm}
    \today
  }

  \vfill
  {\color{edgeblue}\rule{\linewidth}{2pt}}
\end{titlepage}

% --- Tabla de contenidos ---
\tableofcontents
\newpage

% ============================================================
\section{¿Qué es EdgeLeak AI?}
% ============================================================

EdgeLeak AI es una aplicación móvil desarrollada en Flutter que actúa como
\textbf{cerebro central} de un sistema IoT para la detección inteligente de
microfugas en tuberías domésticas. El sistema combina datos de un sensor de
sonido y un sensor de flujo de agua (ESP32) con un modelo de Inteligencia
Artificial (Groq API) para emitir veredictos sobre el estado de la tubería.

Cada vez que la IA detecta una posible fuga, el resultado queda guardado
automáticamente en una base de datos local SQLite dentro del dispositivo que
ejecuta la app.

\begin{tcolorbox}[infobox]
  Este documento explica cómo tu aplicación (Flutter + Firebase) puede
  conectarse a EdgeLeak AI para consultar ese historial de alertas en tiempo
  real, sin necesidad de conocer los detalles internos del sistema.
\end{tcolorbox}

% ============================================================
\section{Cómo funciona la conexión}
% ============================================================

EdgeLeak AI levanta un \textbf{pequeño servidor web} dentro del mismo celular
donde corre la app, escuchando en el puerto \textbf{8080}. Ese servidor expone
una URL (llamada \emph{endpoint}) que puedes llamar desde tu aplicación para
obtener los datos.

La comunicación es sencilla:

\begin{enumerate}[leftmargin=*, label=\textbf{\arabic*.}]
  \item Tu app envía una petición HTTP al celular con EdgeLeak AI.
  \item EdgeLeak AI verifica que tengas la llave secreta correcta.
  \item Si todo está bien, te responde con la lista de alertas en formato JSON.
\end{enumerate}

\begin{tcolorbox}[warningbox]
  Ambos dispositivos (el que corre EdgeLeak AI y el que corre tu app) deben
  estar conectados a la \textbf{misma red Wi-Fi local} para que la comunicación
  funcione. No se necesita internet ni servidor externo.
\end{tcolorbox}

% ============================================================
\section{Requisito previo: obtener la IP del dispositivo}
% ============================================================

Antes de hacer cualquier petición necesitas saber la \textbf{dirección IP local}
del celular donde corre EdgeLeak AI. Piensa en la IP como la "dirección de la
casa" del dispositivo en tu red Wi-Fi.

\subsection{Cómo encontrar la IP en Android}

\begin{enumerate}[leftmargin=*, label=\arabic*.)]
  \item Abre \textbf{Ajustes} en el celular con EdgeLeak AI.
  \item Ve a \textbf{Conexiones} $\rightarrow$ \textbf{Wi-Fi}.
  \item Toca el nombre de la red Wi-Fi a la que estás conectado.
  \item Busca el campo que dice \textbf{Dirección IP} o \textbf{IP address}.
\end{enumerate}

Verás algo como: \texttt{192.168.1.45}

La URL base que usarás en todas tus peticiones será:

\begin{lstlisting}[style=httpStyle]
http://192.168.1.45:8080
\end{lstlisting}

\begin{tcolorbox}[infobox]
  \textbf{Tip:} La IP puede cambiar si el router asigna una nueva al dispositivo.
  Para evitar problemas, puedes configurar una IP fija (estática) en los ajustes
  de Wi-Fi del celular con EdgeLeak AI, o guardar la IP en Firebase para que tu
  app siempre sepa a dónde conectarse.
\end{tcolorbox}

% ============================================================
\section{Autenticación: la llave secreta (x-api-key)}
% ============================================================

Para proteger los datos, EdgeLeak AI exige que todas las peticiones incluyan
una \textbf{llave secreta} en el encabezado (header) de la solicitud HTTP. Esto
evita que cualquier persona en la red pueda leer el historial sin permiso.

\subsection{¿Cómo se envía la llave?}

Se incluye en el encabezado HTTP con el nombre \texttt{x-api-key}:

\begin{lstlisting}[style=httpStyle]
x-api-key: edgeleak-share-2024
\end{lstlisting}

\subsection{Valor de la llave}

\begin{center}
  \begin{tcolorbox}[
    enhanced,
    colback=codebg,
    colframe=edgegray,
    width=0.6\textwidth,
    center,
    left=10pt, right=10pt,
  ]
    \centering
    \texttt{\large edgeleak-share-2024}
  \end{tcolorbox}
\end{center}

\begin{tcolorbox}[warningbox]
  Este valor puede ser personalizado en el archivo de configuración
  \texttt{.env} de EdgeLeak AI (variable \texttt{LOCAL\_API\_KEY}). Si el
  equipo de EdgeLeak cambia la clave, deberás actualizar el valor en tu app
  también. Coordina este cambio con el equipo de EdgeLeak AI.
\end{tcolorbox}

% ============================================================
\section{Endpoints disponibles}
% ============================================================

EdgeLeak AI expone dos endpoints. El que te interesa como consumidor externo
es el de historial.

\begin{center}
\begin{tabular}{>{\bfseries}l l l p{5.5cm}}
  \toprule
  Método & Ruta & Protegido & Descripción \\
  \midrule
  \textcolor{successgreen}{GET}  & \texttt{/api/history} & Sí & Devuelve la lista de alertas de fuga detectadas. \\
  \addlinespace
  \textcolor{edgeorange}{POST} & \texttt{/api/sensor}  & Sí & Recibe datos del sensor ESP32 (uso interno del hardware). \\
  \bottomrule
\end{tabular}
\end{center}

\vspace{0.3cm}
\begin{tcolorbox}[infobox]
  Como aplicación consumidora solo necesitas usar \textbf{GET /api/history}.
  El endpoint \texttt{POST /api/sensor} es exclusivo para el hardware ESP32.
\end{tcolorbox}

% ============================================================
\section{Endpoint principal: GET /api/history}
% ============================================================

\subsection{¿Qué hace?}

Devuelve \textbf{todas las alertas de fuga} que EdgeLeak AI ha registrado,
ordenadas de la más reciente a la más antigua. Cada alerta tiene el veredicto
de la IA, la severidad, un mensaje explicativo y la fecha exacta.

\subsection{Cómo construir la petición}

\begin{lstlisting}[style=httpStyle, caption={Petición HTTP completa}]
GET http://[IP_DEL_DISPOSITIVO]:8080/api/history HTTP/1.1
x-api-key: edgeleak-share-2024
\end{lstlisting}

Reemplaza \texttt{[IP\_DEL\_DISPOSITIVO]} con la IP real del celular con
EdgeLeak AI (ver Sección 3).

\subsection{Respuesta exitosa (código 200)}

La API responde con un arreglo JSON. Cada elemento del arreglo representa
una alerta:

\begin{lstlisting}[style=jsonStyle, caption={Ejemplo de respuesta exitosa}]
[
  {
    "id": 5,
    "veredicto": "Fuga Detectada",
    "severidad": "Critica",
    "mensaje": "Caudal anormalmente bajo (0.04 L/min) con ruido elevado
                detectado. Alta probabilidad de microfuga activa.",
    "timestamp": "2025-04-28T14:32:10.123"
  },
  {
    "id": 4,
    "veredicto": "Sin Fuga",
    "severidad": "Normal",
    "mensaje": "Patron de flujo estable. No se detectan anomalias.",
    "timestamp": "2025-04-28T13:15:42.000"
  },
  {
    "id": 3,
    "veredicto": "Fuga Detectada",
    "severidad": "Alta",
    "mensaje": "Posible fuga de baja intensidad. Se recomienda revision.",
    "timestamp": "2025-04-27T09:08:31.500"
  }
]
\end{lstlisting}

\subsection{Significado de cada campo}

\begin{center}
\begin{tabular}{>{\ttfamily}l l p{6.5cm}}
  \toprule
  \textrm{Campo} & Tipo & Descripción \\
  \midrule
  id          & Número entero   & Identificador único del registro. Se incrementa
                                   automáticamente. \\
  \addlinespace
  veredicto   & Texto           & Resultado del análisis de IA. Puede ser
                                   \textbf{"Fuga Detectada"} o
                                   \textbf{"Sin Fuga"}. \\
  \addlinespace
  severidad   & Texto           & Nivel de gravedad asignado por la IA. Valores
                                   posibles: \textbf{Crítica}, \textbf{Alta},
                                   \textbf{Media}, \textbf{Baja}, \textbf{Normal}. \\
  \addlinespace
  mensaje     & Texto           & Descripción detallada generada por la IA
                                   explicando el resultado. \\
  \addlinespace
  timestamp   & Texto (ISO 8601)& Fecha y hora exacta en que se registró la
                                   alerta. Formato: \texttt{YYYY-MM-DDTHH:mm:ss.SSS} \\
  \bottomrule
\end{tabular}
\end{center}

\subsection{Respuestas de error}

\begin{tcolorbox}[errorbox, title={\faTimesCircle\ Error 401 — No autorizado}]
\begin{lstlisting}[style=jsonStyle]
{
  "error": "No autorizado. Header x-api-key invalido o ausente."
}
\end{lstlisting}
  \vspace{0.2cm}
  \textbf{Causa:} No enviaste el header \texttt{x-api-key} o el valor es incorrecto.\\
  \textbf{Solución:} Verifica que el header esté incluido y que el valor sea exactamente \texttt{edgeleak-share-2024}.
\end{tcolorbox}

\begin{tcolorbox}[errorbox, title={\faTimesCircle\ Error 404 — Endpoint no encontrado}]
\begin{lstlisting}[style=jsonStyle]
{
  "error": "Endpoint no encontrado: GET /api/historial"
}
\end{lstlisting}
  \vspace{0.2cm}
  \textbf{Causa:} La ruta está mal escrita. El endpoint correcto es \texttt{/api/history} (en inglés).\\
  \textbf{Solución:} Corrige la URL.
\end{tcolorbox}

% ============================================================
\section{Integración con Flutter + Firebase}
% ============================================================

A continuación encontrarás el código completo y listo para usar en tu
aplicación Flutter. Está organizado en pasos para que sea fácil de seguir.

\subsection{Paso 1 — Agregar la dependencia HTTP}

En tu archivo \texttt{pubspec.yaml}, asegúrate de tener el paquete \texttt{http}:

\begin{lstlisting}[style=dartStyle, caption={pubspec.yaml}]
dependencies:
  flutter:
    sdk: flutter
  http: ^1.2.0          # Paquete para hacer peticiones HTTP
  firebase_core: ^2.x.x # Ya lo tienes en tu proyecto
\end{lstlisting}

Luego ejecuta en la terminal:

\begin{lstlisting}[style=httpStyle]
flutter pub get
\end{lstlisting}

\subsection{Paso 2 — Crear el modelo de datos}

Crea el archivo \texttt{lib/models/alerta\_fuga.dart} para representar
cada alerta que recibes de EdgeLeak AI:

\begin{lstlisting}[style=dartStyle, caption={lib/models/alerta\_fuga.dart}]
class AlertaFuga {
  final int id;
  final String veredicto;
  final String severidad;
  final String mensaje;
  final DateTime timestamp;

  AlertaFuga({
    required this.id,
    required this.veredicto,
    required this.severidad,
    required this.mensaje,
    required this.timestamp,
  });

  // Convierte el JSON que llega de la API a un objeto AlertaFuga
  factory AlertaFuga.fromJson(Map<String, dynamic> json) {
    return AlertaFuga(
      id: json['id'] as int,
      veredicto: json['veredicto'] as String,
      severidad: json['severidad'] as String,
      mensaje: json['mensaje'] as String,
      timestamp: DateTime.parse(json['timestamp'] as String),
    );
  }

  // Dice si es una fuga detectada (util para mostrar alertas)
  bool get esFuga => veredicto == 'Fuga Detectada';

  // Dice si la severidad es critica o alta
  bool get esGrave => severidad == 'Critica' || severidad == 'Alta';
}
\end{lstlisting}

\subsection{Paso 3 — Crear el servicio de conexión}

Crea el archivo \texttt{lib/services/edgeleak\_service.dart}. Este archivo
es el encargado de hablar con EdgeLeak AI:

\begin{lstlisting}[style=dartStyle, caption={lib/services/edgeleak\_service.dart}]
import 'dart:convert';
import 'package:http/http.dart' as http;
import '../models/alerta_fuga.dart';

class EdgeLeakService {
  // --- CONFIGURA ESTOS VALORES ---
  // Pon aqui la IP del celular donde corre EdgeLeak AI.
  // Si la guardas en Firebase Remote Config, puedes cambiarla
  // sin actualizar la app.
  static const String _ipEdgeLeak = '192.168.1.45';
  static const int    _puerto     = 8080;
  static const String _apiKey     = 'edgeleak-share-2024';
  // --------------------------------

  static String get _baseUrl => 'http://$_ipEdgeLeak:$_puerto';

  /// Obtiene el historial completo de alertas de EdgeLeak AI.
  /// Devuelve una lista de [AlertaFuga] o lanza una excepcion si algo falla.
  static Future<List<AlertaFuga>> obtenerHistorial() async {
    final url = Uri.parse('$_baseUrl/api/history');

    final response = await http.get(
      url,
      headers: {
        'x-api-key': _apiKey,        // La llave secreta
        'Content-Type': 'application/json',
      },
    ).timeout(
      const Duration(seconds: 10),  // Maximo 10 segundos de espera
    );

    if (response.statusCode == 200) {
      // Convierte el texto JSON en una lista de objetos Dart
      final List<dynamic> jsonList = jsonDecode(response.body);
      return jsonList
          .map((item) => AlertaFuga.fromJson(item as Map<String, dynamic>))
          .toList();
    } else if (response.statusCode == 401) {
      throw Exception('API Key incorrecta. Verifica el valor de _apiKey.');
    } else {
      throw Exception(
        'Error del servidor: ${response.statusCode} — ${response.body}',
      );
    }
  }
}
\end{lstlisting}

\subsection{Paso 4 — Usar el servicio en una pantalla}

Crea o edita una pantalla para mostrar el historial. Aquí un ejemplo
completo con \texttt{FutureBuilder}:

\begin{lstlisting}[style=dartStyle, caption={lib/screens/historial\_edgeleak\_screen.dart}]
import 'package:flutter/material.dart';
import '../services/edgeleak_service.dart';
import '../models/alerta_fuga.dart';

class HistorialEdgeLeakScreen extends StatefulWidget {
  const HistorialEdgeLeakScreen({super.key});

  @override
  State<HistorialEdgeLeakScreen> createState() =>
      _HistorialEdgeLeakScreenState();
}

class _HistorialEdgeLeakScreenState
    extends State<HistorialEdgeLeakScreen> {

  // Future que guarda el resultado de la peticion a la API
  late Future<List<AlertaFuga>> _futureHistorial;

  @override
  void initState() {
    super.initState();
    _cargarHistorial();
  }

  void _cargarHistorial() {
    setState(() {
      _futureHistorial = EdgeLeakService.obtenerHistorial();
    });
  }

  @override
  Widget build(BuildContext context) {
    return Scaffold(
      appBar: AppBar(
        title: const Text('Historial EdgeLeak AI'),
        actions: [
          // Boton para refrescar la lista
          IconButton(
            icon: const Icon(Icons.refresh),
            onPressed: _cargarHistorial,
          ),
        ],
      ),
      body: FutureBuilder<List<AlertaFuga>>(
        future: _futureHistorial,
        builder: (context, snapshot) {

          // Mientras espera la respuesta, muestra un indicador
          if (snapshot.connectionState == ConnectionState.waiting) {
            return const Center(child: CircularProgressIndicator());
          }

          // Si hubo un error, lo muestra con opcion de reintentar
          if (snapshot.hasError) {
            return Center(
              child: Column(
                mainAxisSize: MainAxisSize.min,
                children: [
                  const Icon(Icons.wifi_off, size: 64, color: Colors.grey),
                  const SizedBox(height: 12),
                  Text(
                    'No se pudo conectar a EdgeLeak AI.\n'
                    'Verifica que ambos dispositivos esten\n'
                    'en la misma red Wi-Fi.',
                    textAlign: TextAlign.center,
                    style: const TextStyle(color: Colors.grey),
                  ),
                  const SizedBox(height: 16),
                  ElevatedButton.icon(
                    icon: const Icon(Icons.refresh),
                    label: const Text('Reintentar'),
                    onPressed: _cargarHistorial,
                  ),
                ],
              ),
            );
          }

          final alertas = snapshot.data ?? [];

          if (alertas.isEmpty) {
            return const Center(
              child: Text('No hay alertas registradas aun.'),
            );
          }

          // Muestra la lista de alertas
          return ListView.builder(
            itemCount: alertas.length,
            padding: const EdgeInsets.all(8),
            itemBuilder: (context, index) {
              final alerta = alertas[index];
              return _tarjetaAlerta(alerta);
            },
          );
        },
      ),
    );
  }

  Widget _tarjetaAlerta(AlertaFuga alerta) {
    // Color segun severidad
    final Color color = switch (alerta.severidad) {
      'Critica' => Colors.red,
      'Alta'    => Colors.orange,
      'Media'   => Colors.amber,
      _         => Colors.green,
    };

    return Card(
      margin: const EdgeInsets.symmetric(vertical: 6, horizontal: 4),
      child: ListTile(
        leading: Icon(
          alerta.esFuga ? Icons.water_drop : Icons.check_circle,
          color: color,
          size: 36,
        ),
        title: Text(
          alerta.veredicto,
          style: TextStyle(fontWeight: FontWeight.bold, color: color),
        ),
        subtitle: Column(
          crossAxisAlignment: CrossAxisAlignment.start,
          children: [
            Text('Severidad: ${alerta.severidad}'),
            Text(alerta.mensaje, maxLines: 2, overflow: TextOverflow.ellipsis),
            Text(
              '${alerta.timestamp.day}/${alerta.timestamp.month}/'
              '${alerta.timestamp.year}  '
              '${alerta.timestamp.hour}:${alerta.timestamp.minute.toString()
              .padLeft(2, "0")}',
              style: const TextStyle(fontSize: 11, color: Colors.grey),
            ),
          ],
        ),
        isThreeLine: true,
      ),
    );
  }
}
\end{lstlisting}

\subsection{Paso 5 — Guardar en Firebase (opcional pero recomendado)}

Si quieres sincronizar los datos de EdgeLeak AI con tu Firebase para
tenerlos disponibles aunque los dispositivos no estén en la misma red,
puedes hacer esto después de obtener el historial:

\begin{lstlisting}[style=dartStyle, caption={Sincronización con Firebase Firestore}]
import 'package:cloud_firestore/cloud_firestore.dart';

// Dentro de tu clase o controlador:
Future<void> sincronizarConFirebase() async {
  final alertas = await EdgeLeakService.obtenerHistorial();
  final db = FirebaseFirestore.instance;
  final coleccion = db.collection('historial_edgeleak');

  // Sube cada alerta a Firestore usando el id como documento
  for (final alerta in alertas) {
    await coleccion.doc(alerta.id.toString()).set({
      'veredicto':  alerta.veredicto,
      'severidad':  alerta.severidad,
      'mensaje':    alerta.mensaje,
      'timestamp':  alerta.timestamp.toIso8601String(),
      'esFuga':     alerta.esFuga,
      'sincronizadoEn': FieldValue.serverTimestamp(),
    });
  }
}
\end{lstlisting}

\begin{tcolorbox}[infobox, title={\faLightbulb\ Tip — IP dinámica con Firebase Remote Config}]
  Guarda la IP de EdgeLeak AI como un parámetro en Firebase Remote Config
  (ej. \texttt{edgeleak\_ip = "192.168.1.45"}). Así, cuando la IP cambie, el
  administrador la actualiza en Firebase y tu app la lee automáticamente sin
  necesidad de publicar una nueva versión.

  \begin{lstlisting}[style=dartStyle]
final remoteConfig = FirebaseRemoteConfig.instance;
await remoteConfig.fetchAndActivate();
final ip = remoteConfig.getString('edgeleak_ip');
  \end{lstlisting}
\end{tcolorbox}

% ============================================================
\section{Prueba rápida antes de codificar}
% ============================================================

Antes de escribir código, puedes verificar que la conexión funciona usando
\textbf{Postman} o la herramienta \textbf{curl} desde una terminal:

\subsection{Con Postman}

\begin{enumerate}[leftmargin=*, label=\arabic*.)]
  \item Crea una nueva petición de tipo \textbf{GET}.
  \item En la URL escribe: \texttt{http://192.168.1.45:8080/api/history}
        (cambia la IP según corresponda).
  \item Ve a la pestaña \textbf{Headers} y agrega:
        \begin{itemize}
          \item Key: \texttt{x-api-key}
          \item Value: \texttt{edgeleak-share-2024}
        \end{itemize}
  \item Haz clic en \textbf{Send}.
  \item Deberías ver la respuesta JSON con la lista de alertas.
\end{enumerate}

\subsection{Con curl (terminal)}

\begin{lstlisting}[style=httpStyle, caption={Comando curl de prueba}]
curl -X GET "http://192.168.1.45:8080/api/history" \
     -H "x-api-key: edgeleak-share-2024"
\end{lstlisting}

% ============================================================
\section{Resumen rápido de referencia}
% ============================================================

\begin{center}
\begin{tabular}{l l}
  \toprule
  \textbf{Parámetro}       & \textbf{Valor} \\
  \midrule
  Protocolo              & HTTP/1.1 \\
  Puerto                 & 8080 \\
  Endpoint historial     & \texttt{GET /api/history} \\
  Header de autenticación& \texttt{x-api-key} \\
  Valor de la llave      & \texttt{edgeleak-share-2024} \\
  Formato de respuesta   & JSON (arreglo de objetos) \\
  Timeout recomendado    & 10 segundos \\
  Red requerida          & Wi-Fi local compartida \\
  \bottomrule
\end{tabular}
\end{center}

\vspace{0.5cm}

\begin{tcolorbox}[successbox, title={\faCheckCircle\ Lista de verificación antes del primer intento}]
  \begin{itemize}[leftmargin=*, label=\checkmark]
    \item Ambos dispositivos están conectados a la misma red Wi-Fi.
    \item Tienes la IP actual del celular con EdgeLeak AI.
    \item El header \texttt{x-api-key: edgeleak-share-2024} está incluido.
    \item La URL termina en \texttt{/api/history} (no \texttt{/api/historial}).
    \item La app de EdgeLeak AI está abierta y en la pantalla del dashboard.
    \item Probaste la conexión con Postman antes de codificar.
  \end{itemize}
\end{tcolorbox}

% ============================================================
\section{Preguntas frecuentes}
% ============================================================

\begin{description}[style=nextline, leftmargin=0pt]

  \item[\faQuestionCircle\ ¿Qué pasa si la app de EdgeLeak AI está cerrada?]
  El servidor HTTP se apaga cuando la app se cierra. Debes tener EdgeLeak AI
  abierta (en la pantalla del dashboard) para que las peticiones funcionen.

  \item[\faQuestionCircle\ ¿Puedo conectarme desde internet (fuera de la red Wi-Fi)?]
  No directamente. El servidor solo escucha en la red local. Para acceso
  remoto necesitarías configurar un túnel o sincronizar los datos en Firebase
  (ver Paso 5 de la Sección 7).

  \item[\faQuestionCircle\ ¿Cuántas alertas devuelve el endpoint?]
  Devuelve \textbf{todas} las alertas almacenadas en el dispositivo, sin límite.
  Si necesitas paginación, coordínalo con el equipo de EdgeLeak AI para una
  versión futura de la API.

  \item[\faQuestionCircle\ ¿Qué significa el campo \texttt{timestamp}?]
  Es la fecha y hora en que se registró la alerta, en formato ISO 8601
  (estándar internacional). Ejemplo: \texttt{2025-04-28T14:32:10.123} significa
  28 de abril de 2025 a las 2:32 PM. Puedes parsearlo con
  \texttt{DateTime.parse()} en Dart.

  \item[\faQuestionCircle\ ¿Qué hago si recibo error 401?]
  Revisa que el header \texttt{x-api-key} esté presente en la petición y que
  el valor sea exactamente \texttt{edgeleak-share-2024} (sin espacios
  adicionales, distingue mayúsculas y minúsculas).

  \item[\faQuestionCircle\ ¿El campo \texttt{id} es único?]
  Sí. Es un número que se incrementa automáticamente en la base de datos de
  EdgeLeak AI. Puedes usarlo como identificador único al guardar en Firestore.

\end{description}

% ============================================================
\section*{Contacto y soporte}
% ============================================================
\addcontentsline{toc}{section}{Contacto y soporte}

Para dudas sobre esta integración, reportar problemas o solicitar nuevas
funcionalidades en la API, comunícate directamente con el equipo de
\textbf{EdgeLeak AI} a través de los canales del Proyecto Integrador (PICUR)
de la Corporación Universitaria Remington.

\vspace{1cm}
\begin{center}
  {\color{edgeblue}\rule{0.5\linewidth}{0.8pt}}\\[0.3cm]
  {\color{edgegray}\small Documento generado para uso del equipo externo de integración.\\
  EdgeLeak AI — PICUR, Corporación Universitaria Remington — \the\year}
\end{center}

\end{document}
